\chapter{AI Agents -- Konzepte, Frameworks und Orchestrierung}
\label{chap:ai_agents}

% ============================================================
% LERNZIELE
% ============================================================
\begin{lernziele}
Nach diesem Kapitel kannst du:
\begin{itemize}
    \item Den fundamentalen Agent-Loop (Observe-Plan-Act-Reflect) erklären und implementieren
    \item Function Calling / Tool Use für OpenAI, Anthropic und Google produktiv einsetzen
    \item ReAct-Agenten für Single- und Multi-Step-Aufgaben bauen
    \item Multi-Agent-Systeme mit LangGraph orchestrieren (Supervisor + Worker)
    \item Agent-Sicherheit (Sandboxing, Human-in-the-Loop, Token-Budget) produktiv umsetzen
    \item Die Trade-offs zwischen den großen Agent-Frameworks bewerten
\end{itemize}
\end{lernziele}

% ============================================================
% MOTIVATION
% ============================================================
\section{Motivation}

Ein einfacher LLM-Aufruf ist reaktiv: Prompt rein, Antwort raus, fertig. Für viele Anwendungen reicht das -- Klassifikation, Extraktion, Zusammenfassung. Aber sobald ein System \textbf{mehrere Schritte} ausführen muss, \textbf{Tools aufrufen} und \textbf{Ergebnisse interpretieren} soll, reicht ein einzelner Prompt nicht mehr.

Ein AI Agent ist genau das: Ein System, das LLMs als \textbf{Reasoning Engine} nutzt, um autonom Aufgaben zu planen und auszuführen. Der Agent entscheidet selbst, welche Tools er in welcher Reihenfolge aufruft, bewertet Zwischenergebnisse und passt seine Strategie an.

2026 sind AI Agents der dominierende Trend in der LLM-Entwicklung. Während 2023/2024 reine RAG-Systeme dominierten, setzen Produktionsteams heute auf agentische Systeme, die Recherche, Analyse, Entscheidung und Aktion in einem automatisierten Workflow vereinen.

\begin{verbatim}
   Prompt:        [User] --> [LLM] --> [Antwort]

   RAG:           [User] --> [Retriever] --> [LLM] --> [Antwort mit Kontext]

   Agent:         [User] --> [LLM plant] --> [Tool 1] --> [LLM bewertet]
                               |                            |
                               +------- [Tool 2] <----------+
                                            |
                               +------- [Tool 3] <----------+
                               |                            |
                               +--> [Finale Antwort] <------+
\end{verbatim}

\section{Vom Prompt zum System -- Das Denken in Loops}

\subsection{Der Bruch mit dem Prompt-Denken}

Das klassische Prompt-Denken basiert auf Kontrolle im Moment. Du formulierst präzise, um das bestmögliche Ergebnis in einem einzelnen Schritt zu erhalten. Die Qualität hängt stark von deiner Fähigkeit ab, die Anfrage perfekt zu beschreiben. Das Problem: Jede neue Situation erfordert erneut Aufmerksamkeit, erneut Kontext, erneut Entscheidung.

Loops verschieben dieses Modell fundamental. Statt jede Entscheidung selbst zu treffen, definierst du einen Rahmen, innerhalb dessen Entscheidungen automatisch entstehen. Du steuerst nicht mehr den Output. Du steuerst das System, das den Output erzeugt.

\begin{verbatim}
   Prompt-Denken:    [Trigger] --> [LLM] --> [Antwort] --> Ende

   Loop-Denken:      [Trigger] --> [Plan] --> [Act] --> [Evaluate]
                                         ^                  |
                                         |   wiederholen    |
                                         +------------------+
                                                             |
                                         +--> [Finales Ergebnis]
\end{verbatim}

\subsection{Die Loop-Architektur}

Ein funktionierender Loop besteht aus mehreren Ebenen:

\begin{enumerate}[leftmargin=*]
    \item \textbf{Trigger}: Ein zeitbasierter, ereignisgetriebener oder datenbasierter Start. Kein manuelles Eingreifen nötig.
    \item \textbf{Planung}: Das Modell oder ein spezialisierter Planner erzeugt einen Ausführungsplan.
    \item \textbf{Aktion}: Tools, APIs oder Sub-Agenten führen die geplanten Schritte aus.
    \item \textbf{Bewertung}: Das Ergebnis wird gegen objektive Kriterien geprüft (Tests, Builds, Metriken).
    \item \textbf{Persistenz}: Der Zustand wird gespeichert -- ohne Gedächtnis ist jeder Loop nur ein isolierter Versuch.
\end{enumerate}

\subsection{Gedächtnis als Kontinuitätsfaktor}

Der entscheidende Unterschied zwischen einem Prompt und einem System ist nicht die Intelligenz des Modells, sondern seine Fähigkeit, sich über Zeit zu organisieren. Dieses Gedächtnis ist oft unspektakulär implementiert: eine Datei, ein Log, ein strukturiertes Dokument.

\begin{lstlisting}[language=Python, caption=Loop-Gedächtnis als Datei-Persistenz]
class LoopMemory:
    """Ein strukturell einfaches, aber wirkungsvolles Gedaechtnis."""
    def __init__(self, path: str = "loop_state.json"):
        self.path = path
        self.history: list[dict] = []
        self._load()

    def _load(self):
        try:
            with open(self.path) as f:
                data = json.load(f)
                self.history = data.get("history", [])
        except FileNotFoundError:
            self.history = []

    def save(self):
        with open(self.path, "w") as f:
            json.dump({"history": self.history[-100:]}, f)

    def add_step(self, step_type: str, content: str, result: str):
        self.history.append({
            "step": len(self.history) + 1,
            "type": step_type,
            "content": content,
            "result": result,
            "timestamp": time.time(),
        })
        self.save()

    def context(self, last_n: int = 5) -> str:
        recent = self.history[-last_n:]
        return "\n".join(
            f"#{s['step']} [{s['type']}]: {s['content']} -> {s['result'][:200]}"
            for s in recent
        )
\end{lstlisting}

Ohne Gedächtnis ist ein Loop nur ein wiederholter Prompt. Mit Gedächtnis wird er zu einem Prozess, der sich selbst verbessert.

\subsection{Trennung von Erstellung und Bewertung}

Einer der häufigsten Fehler in frühen Agentensystemen ist die Vermischung von Erstellung und Bewertung. Ein System, das seine eigene Arbeit bewertet, neigt zur Selbstbestätigung. Es akzeptiert Ergebnisse, weil sie \glqq gut genug aussehen\grqq, nicht weil sie objektiv korrekt sind.

Deshalb braucht ein stabiler Loop mindestens zwei Perspektiven:

\begin{description}[leftmargin=*,style=nextline]
    \item[Generator] Erzeugt den Output (Code, Text, Entscheidung). Optimiert auf Kreativität und Geschwindigkeit.
    \item[Evaluator] Prüft den Output gegen externe Kriterien (Tests, Builds, Regeln, messbare Bedingungen). Optimiert auf Korrektheit.
\end{description}

\begin{lstlisting}[language=Python, caption=Generator-Evaluator-Trennung im Loop]
class LoopEngine:
    """Trennt Erzeugung und Bewertung in zwei unabhaengige Phasen."""

    def __init__(self, generator, evaluator, max_iterations: int = 5):
        self.generator = generator
        self.evaluator = evaluator
        self.max_iterations = max_iterations
        self.memory = LoopMemory()

    def run(self, task: str) -> str:
        for iteration in range(1, self.max_iterations + 1):
            # Phase 1: Erzeugen
            output = self.generator.generate(task, self.memory.context())
            self.memory.add_step("generate", task, output)

            # Phase 2: Bewerten (durch unabhaengige Instanz)
            eval_result = self.evaluator.evaluate(task, output)
            self.memory.add_step("evaluate", output, str(eval_result))

            if eval_result["passed"]:
                return output

            # Phase 3: Feedback einarbeiten
            task = f"{task}\n\nFeedback: {eval_result['feedback']}"

        raise RuntimeError(f"Nach {self.max_iterations} Iterationen kein Ergebnis.")
\end{lstlisting}

\subsection{Parallelität durch Isolation}

Sobald mehrere Prozesse gleichzeitig arbeiten, entsteht ein neues Problem: Konflikt. Wenn mehrere Agenten denselben Raum verändern, zerstören sie sich gegenseitig. Loops benötigen Isolation -- getrennte Arbeitsumgebungen, getrennte Zustände, getrennte Verantwortung.

\begin{itemize}[leftmargin=*]
    \item \textbf{Dateisystem}: Jeder Loop bekommt ein eigenes Arbeitsverzeichnis
    \item \textbf{Datenbank}: Jeder Loop schreibt in einen eigenen Namespace oder Prefix
    \item \textbf{Kontext}: Jeder Loop hat eine eigene Chat-History, kein geteilter State
    \item \textbf{API-Keys}: Jeder Loop nutzt eigene Rate-Limit-Buckets
\end{itemize}

Parallelität reduziert die Gesamtlaufzeit bei n unabhängigen Tasks von O(n) auf O(1) -- ein entscheidender Faktor für Produktionseinsätze mit hohem Durchsatz.

\subsection{Stop-Bedingungen}

Jeder Loop braucht ein Ende. Ohne definierte Stopp-Bedingung entsteht kein Prozess, sondern eine Schleife ohne Ziel.

Die Stopp-Bedingung muss extern überprüfbar sein:

\begin{itemize}[leftmargin=*]
    \item \textbf{Funktional}: Tests sind bestanden, Build ist erfolgreich, Aufgabe ist abgeschlossen
    \item \textbf{Quantitativ}: Metrik-Schwellwert erreicht (Score $>$ 0.9, Kosten $<$ Budget)
    \item \textbf{Hart}: Maximale Iterationsanzahl überschritten (Safety-Net gegen Endlosschleifen)
\end{itemize}

\begin{verbatim}
   Loop-Abbruchbedingungen:

   [Start] --> [Iteration 1] --> [Test?] --nein--> [Iteration 2] --> ...
                  |                |
                [ja]             [ja]
                  v                v
               [Erfolg]        [Iteration max]
                                  |
                                  v
                              [Fehler/Fallback]
\end{verbatim}

\subsection{Graduierte Autonomie}

Autonomie sollte nicht von null auf hundert gehen, sondern in Stufen wachsen:

\begin{table}[H]
\centering
\begin{tabularx}{\linewidth}{lXl}
\toprule
\textbf{Stufe} & \textbf{Verhalten} & \textbf{Einsatz} \\
\midrule
0 -- Vorschlag & System schlägt vor, Mensch führt aus & Erste Exploration \\
1 -- Entwurf & System erzeugt Entwurf, Mensch prüft und gibt frei & Code-Reviews, Texte \\
2 -- Assistiert & System führt aus, Mensch kann eingreifen (HITL) & Support-Tickets, Datenaufbereitung \\
3 -- Überwacht & System läuft autonom, Mensch prüft Stichproben & Routine-Operationen \\
4 -- Autonom & System läuft vollständig selbstständig & Batch-Jobs, Monitoring \\
\bottomrule
\end{tabularx}
\caption{Stufen der graduierten Autonomie für AI Loops}
\label{tab:autonomie-stufen}
\end{table}

\subsection{Die Verschiebung der Rolle des Entwicklers}

Wenn Loops einmal etabliert sind, verändert sich die Rolle des Entwicklers fundamental. Die tägliche Arbeit besteht nicht mehr darin, einzelne Aufgaben zu lösen, sondern Systeme zu entwerfen, die Aufgaben lösen.

\begin{quote}
Der Fokus verschiebt sich von \glqq Was schreibe ich als Nächstes?\grqq{} zu \glqq Welcher Prozess sollte diese Klasse von Problemen automatisch lösen?\grqq{}
\end{quote}

Der Übergang vom Prompt zum Loop ist kein Feature-Upgrade. Es ist ein Wechsel der Denkweise. Ein Prompt ist ein Gespräch. Ein Loop ist eine Maschine. Sobald diese Maschinen stabil laufen, verändert sich nicht nur die Art, wie Arbeit erledigt wird -- sondern auch, wie Arbeit überhaupt definiert wird.

% ===========================================================%
% GRUNDLAGEN
% ===========================================================%
\section{Grundlagen -- Was ist ein AI Agent?}

\subsection{Der fundamentale Agent-Loop}

Jeder AI Agent folgt einem Loop:

\begin{enumerate}[leftmargin=*]
    \item \textbf{Observe}: Die Eingabe (User-Query, System-Prompt, Kontext) verstehen
    \item \textbf{Plan}: Entscheiden, welcher Schritt als Nächstes kommt
    \item \textbf{Act}: Ein Tool ausführen (API-Call, Datenbank-Query, Berechnung)
    \item \textbf{Reflect}: Das Ergebnis bewerten -- ist das Ziel erreicht? Wenn nein: zurück zu Schritt 2
\end{enumerate}

\begin{verbatim}
   +---------+     +---------+     +---------+     +-----------+
   | Observe | --> |  Plan   | --> |  Act    | --> | Reflect   |
   | Verstehe|     | Entscheide|    | Tool    |    | Bewerte   |
   | Aufgabe |     | naechster|    | ausfueh-|    | Ergebnis  |
   |         |     | Schritt |    | ren     |    |           |
   +---------+     +---------+     +---------+     +-----------+
                       ^                              |
                       |  Ziel nicht erreicht         |
                       +------------------------------+
\end{verbatim}

\subsection{Zentrale Konzepte}

\begin{description}[leftmargin=*,style=nextline]
    \item[Tool] Eine Funktion, die der Agent aufrufen kann: API, Datenbank, Rechner, Dateisystem. Tool haben eine Signatur (Name, Parameter, Rückgabetyp) und eine Beschreibung für das LLM.
    \item[Memory] Der Agent muss sich an frühere Schritte erinnern. Kurzzeitgedächtnis = Chat-History, Langzeitgedächtnis = Vector Store.
    \item[Planning] Der Agent entscheidet, welche Tools in welcher Reihenfolge aufgerufen werden. Bei einfachen Agenten geschieht das implizit (ReAct), bei komplexen explizit (Plan-and-Solve).
    \item[Orchestration] Bei Multi-Agent-Systemen: Wer ruft wen wann auf? Supervisor-Agenten verteilen Aufgaben an spezialisierte Worker.
\end{description}

% ===========================================================%
% ARCHITEKTUR
% ===========================================================%
\section{Architektur -- Agenten-Systeme in Produktion}

\subsection{Agenten-Komponenten}

\begin{verbatim}
   +------------------------------------------------------+
   |                   AI Agent                            |
   |                                                        |
   |  +-----------+  +-----------+  +-------------------+  |
   |  | LLM       |  | Memory    |  | Tool Registry     |  |
   |  | (GPT-4o,  |  | (History, |  | (search, calc,    |  |
   |  |  Claude)  |  |  Vector)  |  |  db_query, ...)   |  |
   |  +-----------+  +-----------+  +-------------------+  |
   |                                                        |
   |  +-----------+  +-----------+  +-------------------+  |
   |  | Planner   |  | Executor  |  | Safety Guard     |  |
   |  | (Strategie)|  | (Tool Run)|  | (HITL, Budget)   |  |
   |  +-----------+  +-----------+  +-------------------+  |
   +------------------------------------------------------+
\end{verbatim}

\subsection{AgentOS / Agent-Betriebssysteme}

AgentOS sind spezialisierte Laufzeitplattformen für autonome Agenten. Sie verwalten nicht nur den Agenten selbst, sondern auch dessen Ressourcen, Berechtigungen, Lebenszyklus, Sicherheit und Überwachung. In einem AgentOS definiert der Entwickler Agenten, Tools, Policies und Workflows; das Betriebssystem führt sie konsistent, robust und sicher aus.

Wichtige Aufgaben eines AgentOS:

\begin{itemize}[leftmargin=*]
    \item \textbf{Lebenszyklus-Management}: Agenten starten, stoppen, neu starten, skalieren und versionieren.
    \item \textbf{Tool-Sandboxing}: Tools und APIs werden isoliert ausgeführt und nur mit den erlaubten Parametern versorgt.
    \item \textbf{Permissions und Policy Enforcement}: Agenten erhalten feingranulare Rechte (z. B. lesender DB-Zugriff, kein Shell-Zugriff) und werden gegen Regeln geprüft.
    \item \textbf{State-Management}: Persistente Speicherbereiche für Conversational Memory, Task-State und Audit-Trails.
    \item \textbf{Orchestrierung}: Multi-Agent-Flows als Graphen, Pipelines oder State Machines.
    \item \textbf{Observability}: Standardisierte Logs, Metriken, Kosten-Tracking, Fehlerraten und Drift.
    \item \textbf{Integrationen}: CI/CD, Monitoring, Ticketing und externe Systeme.
\end{itemize}

\subsubsection{AgentOS-Architektur}

\begin{verbatim}
   +--------------------------------------------------------------+
   |                         AgentOS                              |
   |                                                              |
   |  +----------------+  +----------------+  +------------------+  |
   |  | Agent Runtime  |  | Policy Engine  |  | Observability    |  |
   |  | (LLM, Memory,  |  | (Permissions,  |  | (Logging, Metrics|  |
   |  |  Tool Calls)   |  |  Safety Rules) |  |  + Alerts)       |  |
   |  +----------------+  +----------------+  +------------------+  |
   |                                                              |
   |  +----------------+  +----------------+  +------------------+  |
   |  | Workflow       |  | Tool Sandbox   |  | State Storage    |  |
   |  | Orchestrator   |  | (Connections,  |  | (Memory, Cache,  |  |
   |  |                |  |  Secrets, APIs) |  |  Audit Logs)     |  |
   |  +----------------+  +----------------+  +------------------+  |
   +--------------------------------------------------------------+
\end{verbatim}

AgentOS unterscheiden sich von einem Agent-Framework vor allem durch ihren Betriebsschwerpunkt:

\begin{itemize}[leftmargin=*]
    \item \textbf{Produktionsreife statt Prototyp}: Sie sind für Skalierung, Zuverlässigkeit und Compliance gebaut.
    \item \textbf{Sicherheit als Erstklassiger Bürger}: Policies und Sandboxing sind eingebaut, nicht nachträglich ergänzt.
    \item \textbf{Standardisierte Telemetrie}: Jedes Agenten-Event ist abrufbar, vergleichbar und auditierbar.
    \item \textbf{Plattform-Agnostik}: Agenten laufen über verschiedene LLM-Provider hinweg konsistent.
\end{itemize}

\subsubsection{AgentOS-Prinzipien}

Ein AgentOS ist dann erfolgreich, wenn es diese Prinzipien erfüllt:

\begin{itemize}[leftmargin=*]
    \item \textbf{Konsistente Ausführung}: Der gleiche Agent verhält sich in Test, Staging und Produktion gleich.
    \item \textbf{Sichere Isolation}: Tools, Secrets und externe Verbindungen werden getrennt vom Kernmodell betrieben.
    \item \textbf{Transparente Kontrolle}: Jede Entscheidung, jeder Tool-Call und jeder State-Übergang ist nachvollziehbar.
    \item \textbf{Adaptives Failover}: Agenten können bei Fehlern automatisch zurückfallen, neu starten oder auf einen sicheren Zustand wechseln.
    \item \textbf{Kostenbewusstsein}: Der Betrieb umfasst Budgetkontrolle, Kostenalarme und Modell-Tiering.
\end{itemize}

\subsubsection{Beispiele für AgentOS-Funktionen}

Typische AgentOS-Funktionen, die in einfachen Frameworks fehlen:

\begin{itemize}[leftmargin=*]
    \item \textbf{Task-Queues}: Von Agenten erzeugte Aufgaben werden verwaltet und priorisiert.
    \item \textbf{Retry-Strategien}: Fehlgeschlagene Aktionen werden kontrolliert erneut versucht.
    \item \textbf{Circuit Breaker}: Bei zu vielen Fehlern wird ein Tool oder eine Verbindung temporär deaktiviert.
    \item \textbf{Quota-Management}: Token-, API- und Kosten-Limits werden pro Agent oder Team gesetzt.
    \item \textbf{Feature Flags}: Neue Agentenfähigkeiten werden schrittweise ausgerollt.
\end{itemize}

\subsubsection{AgentOS in der Praxis}

In der Praxis wird die AgentOS-Schicht oft als \textbf{Agent Runtime}, \textbf{Agent Execution Layer} oder \textbf{AI Orchestrator} bezeichnet. Wichtige Beispiele aus der Industrie sind:

\begin{itemize}[leftmargin=*]
    \item \textbf{OpenAI Assistants Platform}: Managed Agent-Betrieb mit Policies, Tool-Integration und Observability.
    \item \textbf{AutoGen / Microsoft}: Orchestrierung von Agenten und Workflows über definierte Task-Messages.
    \item \textbf{LangGraph}: Graphenbasierte State-Machine-Ausführung für Multi-Agent-Flows.
    \item \textbf{CrewAI}: Rollenbasierte Multi-Agenten, die spezialisierte Aufgaben koordinieren.
\end{itemize}

\subsubsection{AgentOS vs. Agent-Framework}

Der Unterschied lässt sich so zusammenfassen:

\begin{description}[leftmargin=*,style=nextline]
    \item[Agent-Framework] erlaubt das Bauen von Agenten: Tool-Calls, Loops, Planung.
    \item[AgentOS] betreibt Agenten: Lebenszyklus, Sicherheit, Monitoring, Compliance.
\end{description}

Ein Projekt kann mit einem Framework starten, aber sobald Agenten in Produktion gehen, braucht es einen AgentOS-ähnlichen Betrieb.

\subsection{Single-Agent vs. Multi-Agent}

\begin{table}[H]
\centering
\begin{tabularx}{\linewidth}{lXX}
\toprule
\textbf{Aspekt} & \textbf{Single-Agent} & \textbf{Multi-Agent} \\
\midrule
Komplexität & Ein LLM, n Tools & Mehrere spezialisierte LLMs \\
Planung & Implizit (ReAct-Loop) & Explizit (Supervisor delegiert an Worker) \\
Fehleranfälligkeit & Ein Fehler = gesamter Fehler & Teilausfälle durch andere Agenten auffangbar \\
Kosten & Ein LLM-Call pro Schritt & Mehrere LLM-Calls pro Schritt (teurer) \\
Latenz & Niedriger & Höher durch Kommunikation \\
Use Case & Klar definierte Aufgaben & Komplexe, multi-disziplinäre Aufgaben \\
\midrule
\end{tabularx}
\caption{Single-Agent vs. Multi-Agent}
\label{tab:agent-types}
\end{table}

% ===========================================================%
% THEORIE
% ===========================================================%
\section{Theorie -- Warum Agenten funktionieren}

\subsection{Augmented LLMs}

Der Grundgedanke: Ein LLM allein kann keine externen Aktionen ausführen. Es ist ein reiner Textgenerator. Durch die Kopplung mit Tools (\textbf{Augmented LLM}) wird es handlungsfähig:

\begin{verbatim}
   Basismodell:    Text-in -> Text-out

   Augmented LLM:  Text-in -> Tool-Decision -> Tool-Result -> Text-out

   Der entscheidende Unterschied:
   Das LLM gibt nicht nur Text aus, sondern auch die Entscheidung,
   welches Tool aufgerufen werden soll.
\end{verbatim}

\subsection{Tool-Use als API-Design-Problem}

Das LLM versteht Tools durch \textbf{natürlichsprachliche Beschreibungen}. Die Qualität des Tool-Use hängt direkt von der Qualität der Tool-Beschreibungen ab:

\begin{itemize}[leftmargin=*]
    \item \textbf{Name}: Kurz, präzise, deskriptiv ("search\_knowledge\_base", nicht "search" oder "func\_1")
    \item \textbf{Description}: Was tut das Tool? Wann sollte es verwendet werden? Welche Parameter sind kritisch?
    \item \textbf{Parameters}: Typen, Enums, Beschreibungen, Required-Flags
    \item \textbf{Output}: Was gibt das Tool zurück? (Das LLM muss den Output interpretieren können)
\end{itemize}

\subsection{Planung vs. Reaktivität}

Es gibt zwei grundlegende Agent-Strategien:

\begin{itemize}[leftmargin=*]
    \item \textbf{Reaktiv (ReAct)}: Der Agent denkt und handelt in einem Loop. Keine vorausschauende Planung -- er entscheidet von Schritt zu Schritt. Gut für einfache, lineare Aufgaben.
    \item \textbf{Planend (Plan-and-Solve)}: Der Agent erstellt zuerst einen vollständigen Plan, führt ihn dann schrittweise aus und passt ihn bei Bedarf an. Besser für komplexe Aufgaben mit vielen Abhängigkeiten.
\end{itemize}

% ===========================================================%
% FUNCTION CALLING
% ===========================================================%
\section{Function Calling / Tool Use}

Tool Use ist das Fundament aller Agenten.

\subsection{OpenAI Function Calling}

\begin{lstlisting}[language=Python, caption=Function Calling mit OpenAI]
tools = [
    {
        "type": "function",
        "function": {
            "name": "get_weather",
            "description": "Wetter fuer eine Stadt abrufen",
            "parameters": {
                "type": "object",
                "properties": {
                    "city": {"type": "string", "description": "Stadtname"},
                    "country_code": {
                        "type": "string",
                        "enum": ["DE", "AT", "CH"],
                    },
                },
                "required": ["city", "country_code"],
            },
        },
    }
]

response = client.chat.completions.create(
    model="gpt-4o",
    messages=[{"role": "user",
               "content": "Wetter morgen in Berlin?"}],
    tools=tools,
)

tool_call = response.choices[0].message.tool_calls[0]
args = json.loads(tool_call.function.arguments)
result = get_weather(args["city"], args["country_code"])

# Ergebnis an das Model zurueckgeben
response = client.chat.completions.create(
    model="gpt-4o",
    messages=[
        {"role": "user", "content": "Wetter morgen in Berlin?"},
        response.choices[0].message,
        {"role": "tool", "tool_call_id": tool_call.id,
         "content": json.dumps(result)},
    ],
    tools=tools,
)
print(response.choices[0].message.content)
\end{lstlisting}

\subsection{Anthropic Tool Use}

\begin{lstlisting}[language=Python, caption=Tool Use mit Anthropic Claude]
response = client.messages.create(
    model="claude-sonnet-4-20250514",
    max_tokens=1024,
    messages=[{"role": "user", "content": "Wetter in Berlin?"}],
    tools=[{
        "name": "get_weather",
        "description": "Wetterdaten abrufen",
        "input_schema": {
            "type": "object",
            "properties": {
                "city": {"type": "string"},
                "country": {"type": "string"},
            },
            "required": ["city"],
        },
    }],
)

for block in response.content:
    if block.type == "tool_use":
        result = get_weather(**block.input)
        response = client.messages.create(
            model="claude-sonnet-4-20250514",
            messages=response.messages + [{
                "role": "user",
                "content": [{
                    "type": "tool_result",
                    "tool_use_id": block.id,
                    "content": str(result),
                }],
            }],
        )
\end{lstlisting}

\subsection{Multi-Provider Tool Abstraction}

In Produktion willst du nicht pro Provider unterschiedliche Tool-Formate pflegen:

\begin{lstlisting}[language=Python, caption=Abstraktion fuer Multi-Provider Tool Use]
from dataclasses import dataclass, field

@dataclass
class AgentTool:
    name: str
    description: str
    fn: callable
    parameters: dict = field(default_factory=dict)

    def to_openai(self) -> dict:
        return {
            "type": "function",
            "function": {
                "name": self.name,
                "description": self.description,
                "parameters": self.parameters,
            },
        }

    def to_anthropic(self) -> dict:
        return {
            "name": self.name,
            "description": self.description,
            "input_schema": self.parameters,
        }

class ToolRegistry:
    def __init__(self):
        self.tools: dict[str, AgentTool] = {}

    def register(self, tool: AgentTool):
        self.tools[tool.name] = tool

    def execute(self, name: str, **kwargs):
        tool = self.tools.get(name)
        if not tool:
            raise ValueError(f"Unknown tool: {name}")
        return tool.fn(**kwargs)
\end{lstlisting}

% ===========================================================%
% AGENT-PATTERNS
% ===========================================================%
\section{Agent-Patterns}

\subsection{ReAct -- der Standard}

\begin{lstlisting}[language=Python, caption=ReAct Agent]
class ReActAgent:
    def __init__(self, tools: ToolRegistry, llm_client, max_steps: int = 10):
        self.tools = tools
        self.llm = llm_client
        self.max_steps = max_steps

    def run(self, user_query: str) -> str:
        messages = [{"role": "system", "content": (
            "Loese Aufgaben durch Thought -> Action -> Observation. "
            "Thought: dein Gedankengang. "
            "Action: Tool-Name. "
            "Action Input: JSON. "
            "Wenn fertig: Answer: Endergebnis."
        )}, {"role": "user", "content": user_query}]

        for step in range(self.max_steps):
            response = self.llm.chat.completions.create(
                model="gpt-4o",
                messages=messages,
                tools=[t.to_openai() for t in self.tools.tools.values()],
            )
            msg = response.choices[0].message

            if msg.content and "Answer:" in msg.content:
                return msg.content

            if msg.tool_calls:
                for tc in msg.tool_calls:
                    args = json.loads(tc.function.arguments)
                    result = self.tools.execute(tc.function.name, **args)
                    messages.append(msg)
                    messages.append({
                        "role": "tool",
                        "tool_call_id": tc.id,
                        "content": str(result),
                    })

        return "Max steps reached without completion."
\end{lstlisting}

\subsection{Plan-and-Solve}

Bei komplexen Aufgaben plant der Agent zuerst den gesamten Ablauf:

\begin{lstlisting}[language=Python, caption=Plan-and-Solve Agent]
class PlanAndSolveAgent:
    def run(self, task: str) -> str:
        # Phase 1: Plan erstellen
        plan_response = self.llm.chat.completions.create(
            model="gpt-4o",
            messages=[{"role": "user", "content": (
                f"Erstelle einen detaillierten Plan fuer diese Aufgabe. "
                f"Liste die Schritte 1..N auf:\n{task}"
            )}],
        )
        plan = plan_response.choices[0].message.content
        steps = self._parse_steps(plan)

        # Phase 2: Schritte ausfuehren
        results = []
        for step in steps:
            result = self._execute_step(step)
            results.append(result)

        # Phase 3: Zusammenfassen
        return self._synthesize(task, results)

    def _parse_steps(self, plan: str) -> list[str]:
        return [line.strip() for line in plan.split("\n")
                if line.strip().startswith(("1.", "2.", "3.", "4.", "5."))]
\end{lstlisting}

\subsection{Multi-Agent orchestrieren}

Supervisor delegiert an spezialisierte Worker:

\begin{lstlisting}[language=Python, caption=Supervisor + Worker Pattern]
class SupervisorAgent:
    def __init__(self):
        self.workers = {
            "researcher": ResearcherAgent(),
            "analyst": AnalystAgent(),
            "writer": WriterAgent(),
        }

    def run(self, task: str) -> str:
        # Schritt 1: Entscheiden, welche Worker gebraucht werden
        plan = self._plan(task)
        results = {}

        # Schritt 2: Worker in der richtigen Reihenfolge ausfuehren
        for worker_name in plan["worker_sequence"]:
            worker = self.workers[worker_name]
            results[worker_name] = worker.run(
                task=plan["tasks"][worker_name],
                context=results,
            )

        # Schritt 3: Ergebnisse zusammenfuehren
        return self._synthesize(results)
\end{lstlisting}

\subsection{LangGraph: Graph-basierte Orchestrierung}

\begin{lstlisting}[language=Python, caption=Multi-Agent-Workflow mit LangGraph]
from langgraph.graph import StateGraph, MessagesState, END

class AgentWorkflow:
    def __init__(self):
        graph = StateGraph(MessagesState)

        graph.add_node("supervisor", self._supervisor)
        graph.add_node("researcher", self._researcher)
        graph.add_node("analyst", self._analyst)
        graph.add_node("writer", self._writer)

        graph.set_entry_point("supervisor")
        graph.add_conditional_edges(
            "supervisor",
            self._route,
            {"research": "researcher", "end": END},
        )
        graph.add_edge("researcher", "analyst")
        graph.add_edge("analyst", "writer")
        graph.add_edge("writer", END)

        self.app = graph.compile()

    async def run(self, user_input: str) -> str:
        result = await self.app.ainvoke(
            {"messages": [("user", user_input)]}
        )
        return result["messages"][-1].content
\end{lstlisting}

% ===========================================================%
% PRAXIS
% ===========================================================%
\section{Praxisbeispiele aus der Industrie}

\subsection{Support-Agent mit Tool-Use}

Ein E-Commerce-Unternehmen ersetzt sein Level-1-Support-Team durch einen AI Agenten:

\begin{itemize}[leftmargin=*]
    \item \textbf{Tools}: Bestellstatus-API, Retoure-Logik, FAQ-Vector-Store, Rabatt-Rechner
    \item \textbf{Workflow}: Agent analysiert die Anfrage $\rightarrow$ sucht in der Wissensdatenbank $\rightarrow$ prüft Bestelldaten $\rightarrow$ generiert Antwort
    \item \textbf{Eskalation}: Bei unsicheren Fällen $\rightarrow$ Human-in-the-Loop (Ticket an menschlichen Support)
    \item \textbf{Ergebnis}: 65\,\% der Tickets vollautomatisiert gelöst, durchschnittliche Lösungszeit von 4h auf 2 Minuten gesenkt
\end{itemize}

\subsection{Code-Review-Agent}

Ein SaaS-Unternehmen nutzt einen Multi-Agent-Workflow für automatische Code-Reviews:

\begin{verbatim}
   [PR geoeffnet] -> [Supervisor]
                       |
           +-----------+-----------+
           |           |           |
       [Security]  [Quality]  [Performance]
           |           |           |
           +-----------+-----------+
                       |
                   [Reviewer]
                       |
                   [Feedback im PR]
\end{verbatim}

\subsection{Finance Research Agent}

Ein Research-Agent analysiert Quartalsberichte:

\begin{enumerate}[leftmargin=*]
    \item \textbf{Recherche}: Extrahiert Finanzkennzahlen aus 10-Q/10-K-Berichten
    \item \textbf{Analyse}: Berechnet Kennzahlen (P/E-Ratio, EBITDA-Marge, FCF-Yield)
    \item \textbf{Vergleich}: Vergleicht mit Wettbewerbern aus der Vector DB
    \item \textbf{Bericht}: Generiert strukturierte Analyse mit Quellenangabe
\end{enumerate}

% ===========================================================%
% MULTI-AGENT
% ===========================================================%
\section{Multi-Agent Orchestrierung}
\label{sec:multi-agent}

Ein einzelner Agent stösst an Grenzen: eine Perspektive, ein Tool-Set, eine Fehlertoleranz. \textbf{Multi-Agent-Systeme} lösen das durch Spezialisierung und Koordination.

\subsection{Architekturmuster}

Drei Muster haben sich in der Produktion bewährt:

\begin{verbatim}
   Supervisor-Worker:
   [User] --> [Supervisor] --> [Worker A: Recherche]
                          --> [Worker B: Analyse]
                          --> [Worker C: Bericht]
                               |
                          [Supervisor aggregiert]
                               |
                          [Antwort an User]
\end{verbatim}

\begin{verbatim}
   Pipeline (Chain):
   [User] --> [Agent A: Extraktion] --> [Agent B: Analyse]
               --> [Agent C: Generierung]
\end{verbatim}

\begin{verbatim}
   Debate / Reflection:
   [Agent A: Pro] --+--> [Mediator]
   [Agent B: Contra] -+          |
                             [Synthese]
\end{verbatim}

\subsubsection{Supervisor-Worker (empfohlen)}

Der Supervisor analysiert die Aufgabe, delegiert an spezialisierte Worker, sammelt Ergebnisse und synthetisiert die Antwort.

\begin{lstlisting}[language=Python, caption=Supervisor-Worker-Architektur]
class SupervisorAgent:
    """Koordiniert spezialisierte Worker-Agenten."""

    def __init__(self, workers: dict[str, callable]):
        self.workers = workers
        self.llm = get_llm()

    async def run(self, task: str) -> dict:
        plan = await self.llm.complete(
            f"Analysiere die Aufgabe und erstelle einen Plan "
            f"mit Workern: {list(self.workers.keys())}."
            f"\nAufgabe: {task}"
        )
        results = {}
        for step in self._parse_steps(plan):
            w = step["worker"]
            results[w] = await self.workers[w](step["task"])

        synthesis = await self.llm.complete(
            f"Synthetisiere Ergebnisse zu einer Antwort:\n"
            f"{json.dumps(results, indent=2)}"
        )
        return {"plan": plan, "results": results, "answer": synthesis}

    def _parse_steps(self, plan: str) -> list[dict]:
        steps = []
        for line in plan.split("\n"):
            if "->" in line and ":" in line:
                worker, task = line.split(":", 1)
                steps.append({
                    "worker": worker.strip().split("->")[-1].strip(),
                    "task": task.strip(),
                })
        return steps
\end{lstlisting}

\subsubsection{Agent-to-Agent-Kommunikation}

In komplexeren Systemen kommunizieren Agenten direkt -- über Message Bus, Blackboard oder Direct Messaging.

\begin{lstlisting}[language=Python, caption=Message-Bus-fuer-Agenten]
class AgentMessageBus:
    """Asynchrone Kommunikation zwischen Agenten."""

    def __init__(self):
        self.queues: dict[str, asyncio.Queue] = {}
        self.results: dict[str, list[dict]] = {}

    def register(self, agent_id: str):
        self.queues[agent_id] = asyncio.Queue()
        self.results[agent_id] = []

    async def send(self, sender: str, recipient: str, msg: dict):
        await self.queues[recipient].put({
            "from": sender, **msg,
        })

    async def receive(self, agent_id: str,
                      timeout: float = 30.0) -> dict:
        return await asyncio.wait_for(
            self.queues[agent_id].get(), timeout=timeout
        )

    def publish(self, sender: str, topic: str, data: dict):
        for agent_id in self.queues:
            if agent_id != sender:
                self.queues[agent_id].put_nowait({
                    "from": sender, "topic": topic, **data,
                })
\end{lstlisting}

\subsection{Frameworks im Vergleich}

\begin{table}[H]
\centering
\begin{tabularx}{\linewidth}{lXXX}
\toprule
\textbf{Kriterium} & \textbf{CrewAI} & \textbf{AutoGen} & \textbf{LangGraph} \\
\midrule
Architektur & Rollenbasiert & Conversational & Graph-basiert \\
Kommunikation & Task-basiert & Free-form Chat & Zustandsautomat \\
State-Management & Automatisch & Manuell & Explizit (Graph) \\
Fehler-Toleranz & Retry + Fallback & Exception Handling & Conditional Edges \\
Lernkurve & Niedrig & Mittel & Steil \\
\bottomrule
\end{tabularx}
\caption{Multi-Agent-Frameworks im Vergleich}
\label{tab:multi-agent-frameworks}
\end{table}

\subsection{Fehlermodi in Multi-Agent-Systemen}

Spezifische Fehler, die in Single-Agent-Systemen nicht auftreten:

\begin{enumerate}[leftmargin=*]
    \item \textbf{Agenten-Loop}: A delegiert an B, B an A -- Endlosschleife. \textbf{Gegenmassnahme}: max\_hops, Zirkelreferenz-Erkennung.
    \item \textbf{Context-Vergiftung}: Falsche Zwischenergebnisse werden als Wahrheit weitergereicht. \textbf{Gegenmassnahme}: Jeder Agent validiert Inputs.
    \item \textbf{Supervisor-Flaschenhals}: Bei $>$ 10 Workern wird Supervisor zum Bottleneck. \textbf{Gegenmassnahme}: Hierarchische Supervisoren.
    \item \textbf{Kosten-Explosion}: 5 Agenten $\times$ 3 Schritte $\times$ 4 LLM-Calls = 60 Calls/Task. \textbf{Gegenmassnahme}: Cost-Budget pro Task, Modell-Tiering.
    \item \textbf{Deadlock}: A wartet auf B, B auf A. \textbf{Gegenmassnahme}: Timeout auf jede Kommunikation.
\end{enumerate}

\subsection{Praxis: Single-Agent vs Multi-Agent}

\begin{verbatim}
   Single-Agent: 3-5 Tools, klar definiert, budget-sensitiv
   -> Beispiel: Uebersetzungs-Assistent

   Multi-Agent (Supervisor-Worker):
      Unterschiedliche Expertise, Exploration + Analyse + Report
   -> Beispiel: Marktforschungs-Assistent

   Multi-Agent (Pipeline):
      Klarer Datenfluss, Stufen unabhaengig testbar
   -> Beispiel: Dokumentenverarbeitung
\end{verbatim}

% ===========================================================%
% BEST PRACTICES
% ===========================================================%
\section{Best Practices}

\begin{itemize}[leftmargin=*]
    \item \textbf{Tool-Beschreibungen sind das halbe System}: Investiere Zeit in gute Tool-Descriptions. Ein schlecht beschriebenes Tool wird nie aufgerufen.
    \item \textbf{Tool-Output begrenzen}: Tool-Ergebnisse werden Teil des LLM-Kontexts. Schneide sie auf das Wesentliche zusammen (max 2.000 Tokens pro Tool-Call).
    \item \textbf{max\_steps setzen}: Ein Agent ohne maximale Schrittanzahl kann in einer Endlosschleife landen. 10--15 Schritte sind ein guter Standard.
    \item \textbf{Human-in-the-Loop}: Bei kritischen Aktionen (Zahlungen, Löschungen, Freigaben) muss ein Mensch bestätigen.
    \item \textbf{Agent-Status loggen}: Jeder Schritt (Thought, Action, Observation) wird geloggt. Das ist die Grundlage für Debugging und Evaluierung.
    \item \textbf{Weniger Tools ist mehr}: Maximal 5--10 Tools gleichzeitig. Bei mehr: hierarchische Tool-Auswahl (Tool-Router delegiert an Sub-Tools).
    \item \textbf{Memory begrenzen}: Die Chat-History wächst mit jedem Schritt. Komprimiere sie regelmäßig (siehe Kapitel~\ref{chap:token_management}).
\end{itemize}

% ===========================================================%
% ANTI-PATTERNS
% ===========================================================%
\section{Anti-Patterns}

\begin{itemize}[leftmargin=*]
    \item \textbf{Zu viele Tools}: Je mehr Tools, desto schlechter die Auswahl. Ein Agent mit 20+ Tools wählt oft das falsche.
    \item \textbf{Kein Timeout}: Ein Agent, der in einer Schleife steckt, generiert unbegrenzt Tokens. Setze immer max\_steps und ein Zeit-Limit.
    \item \textbf{Tool-Ergebnisse nicht zusammenfassen}: Ein Tool, das 10.000 Zeilen JSON zurückgibt, frisst den gesamten Kontext. Fasse zusammen.
    \item \textbf{Agent ohne Memory-Begrenzung}: Nach 20 Schritten ist der Kontext voll. Komprimiere oder truncate.
    \item \textbf{LLM als alleinigen Entscheider bei kritischen Aktionen}: Kein LLM sollte allein Zahlungen auslösen oder Daten löschen. Immer HITL.
    \item \textbf{Agenten für einfache Aufgaben}: Ein Agent ist Overkill, wenn ein einzelner Prompt reicht. Nicht jede Aufgabe braucht einen Loop.
    \item \textbf{Parallel-Tool-Calling ignorieren}: OpenAI unterstützt parallele Tool-Calls. Nutze sie -- sie sparen Zeit und Tokens.
\end{itemize}

% ===========================================================%
% PERFORMANCE
% ===========================================================%
\section{Performance und Optimierung}

\subsection{Latenz pro Agent-Schritt}

\begin{table}[H]
\centering
\begin{tabularx}{\linewidth}{lXr}
\toprule
\textbf{Schritt} & \textbf{Dauer} & \textbf{Kosten} \\
\midrule
LLM-Planung (Thought) & 300--800ms & Input-Tokens \\
Tool-Ausführung & 50--500ms (API-abhängig) & Keine LLM-Kosten \\
Tool-Ergebnis verarbeiten & 100--300ms & Output-Tokens \\
Nächster Schritt planen & 300--800ms & Input-Tokens (History) \\
\midrule
\end{tabularx}
\caption{Latenz und Kosten pro Agent-Schritt (GPT-4o)}
\label{tab:agent-latency}
\end{table}

\subsection{Optimierung}

\begin{itemize}[leftmargin=*]
    \item \textbf{Parallel Tool Calling}: OpenAI kann mehrere Tools in einem Response aufrufen. Das reduziert die Anzahl der LLM-Calls um 30--50\,\%.
    \item \textbf{Günstigeres Modell für einfache Schritte}: Für Routine-Schritte (Formatierung, einfache Extraktion) GPT-4o-mini nutzen, nur für komplexe Reasoning-Schritte GPT-4o.
    \item \textbf{Prompt-Caching für System-Prompts}: Der Agent-System-Prompt ändert sich nie -- cache ihn.
    \item \textbf{Chat-History-Kompression}: Nach jedem N Schritten die History komprimieren (siehe Kapitel~\ref{chap:token_management}).
    \item \textbf{Agent-Caching}: Identische Tasks (gleicher Input, gleicher Kontext) werden gecached.
\end{itemize}

% ===========================================================%
% SICHERHEIT
% ===========================================================%
\section{Sicherheit}

\subsection{Agent-spezifische Risiken}

Agenten sind mächtiger als einfache LLM-Calls -- und damit gefährlicher:

\begin{itemize}[leftmargin=*]
    \item \textbf{Unbegrenzte Tool-Ausführung}: Ein Agent kann Tausende API-Calls in einer Minute auslösen (Kosten-Explosion, Denial-of-Wallet)
    \item \textbf{Prompt Injection via Tools}: Ein Tool-Ergebnis kann schädliche Anweisungen enthalten ("Ignore all instructions...")
    \item \textbf{Data Exfiltration}: Der Agent könnte Daten aus internen Systemen extrahieren und via Tools nach außen senden
    \item \textbf{Command Execution}: Ein Agent mit Shell-/Code-Execution-Tool kann schädliche Befehle ausführen
\end{itemize}

\subsection{Gegenmaßnahmen}

\begin{lstlisting}[language=Python, caption=Agent-Sicherheits-Guard]
class SafetyGuard:
    def __init__(self, max_steps: int = 20,
                 max_tokens_per_step: int = 4000,
                 max_cost_per_run: float = 1.0):
        self.max_steps = max_steps
        self.max_tokens = max_tokens_per_step
        self.max_cost = max_cost_per_run
        self.step_count = 0
        self.total_cost = 0.0

    def check_before_tool(self, tool_name: str, args: dict) -> bool:
        self.step_count += 1
        if self.step_count > self.max_steps:
            raise RuntimeError("Max steps exceeded")

        # Blacklist gefaehrlicher Tools
        if tool_name in {"execute_shell", "delete_file", "send_email"}:
            return False  # Human-in-the-Loop erforderlich

        # Whitelist erlaubter Domains
        if tool_name == "http_request":
            url = args.get("url", "")
            if not url.startswith("https://api.company-intern.com"):
                return False

        return True

    def check_after_tool(self, result: str) -> str:
        # Tool-Ergebnis auf schädliche Anweisungen scannen
        injection_patterns = [
            "ignore previous", "ignore all", "system prompt",
        ]
        for pattern in injection_patterns:
            if pattern.lower() in result.lower():
                return "[Safety: Output blocked - injection detected]"

        # Groesse begrenzen
        max_chars = 10000
        return result[:max_chars]
\end{lstlisting}

\subsection{Human-in-the-Loop (HITL)}

Für kritische Aktionen muss ein Mensch bestätigen:

\begin{lstlisting}[language=Python, caption=Human-in-the-Loop Pattern]
class HITLManager:
    def __init__(self):
        self.pending_approvals = {}

    def request_approval(self, action: str, context: dict) -> str:
        approval_id = str(uuid.uuid4())
        self.pending_approvals[approval_id] = {
            "action": action,
            "context": context,
            "status": "pending",
        }
        # Benachrichtigung an Slack / Teams / E-Mail
        notify_human({
            "type": "agent_approval",
            "approval_id": approval_id,
            "action": action,
            "context": context,
        })
        return approval_id

    def approve(self, approval_id: str) -> bool:
        if approval_id in self.pending_approvals:
            self.pending_approvals[approval_id]["status"] = "approved"
            return True
        return False

    def reject(self, approval_id: str) -> bool:
        if approval_id in self.pending_approvals:
            self.pending_approvals[approval_id]["status"] = "rejected"
            return True
        return False

    def wait_for_approval(self, approval_id: str,
                          timeout: int = 300) -> bool:
        start = time.time()
        while time.time() - start < timeout:
            status = self.pending_approvals[approval_id]["status"]
            if status == "approved":
                return True
            if status == "rejected":
                return False
            time.sleep(1)
        return False  # Timeout
\end{lstlisting}

% ===========================================================%
% ZUSAMMENFASSUNG
% ===========================================================%
\section{Zusammenfassung}

\begin{itemize}[leftmargin=*]
    \item Ein AI Agent folgt einem Observe-Plan-Act-Reflect-Loop -- autonom, multi-step, tool-gestützt
    \item Function Calling / Tool Use ist die Basis: Das LLM entscheidet, welches Tool aufgerufen wird
    \item ReAct ist das Standard-Pattern: Thought $\rightarrow$ Action $\rightarrow$ Observation $\rightarrow$ Wiederholen
    \item Multi-Agent-Systeme (Supervisor + Worker) skalieren auf komplexe Aufgaben
    \item Tool-Beschreibungen sind die halbe Miete -- je besser beschrieben, desto besser der Agent
    \item max\_steps, Token-Budget und Human-in-the-Loop sind Sicherheits-Grundvoraussetzungen
    \item Agenten sind mächtig, aber teuer und langsam im Vergleich zu einzelnen Prompts
\end{itemize}

% ===========================================================%
% MERKE-BOX
% ===========================================================%
\begin{merke}
\begin{itemize}
    \item \textbf{Agent ist nicht immer besser}: Ein Prompt reicht oft. Agent nur bei Multi-Step-Aufgaben.
    \item \textbf{Tools sind die Superkraft}: Jedes Tool braucht eine exzellente Description.
    \item \textbf{max\_steps ist Pflicht}: Ohne Begrenzung läuft der Agent unendlich.
    \item \textbf{HITL für kritische Aktionen}: Kein LLM sollte allein löschen, zahlen oder freigeben.
    \item \textbf{Parallel Tool Calling spart 30--50\,\% der LLM-Calls}.
\end{itemize}
\end{merke}

% ===========================================================%
% PRAXISPROJEKT
% ===========================================================%
\section{Praxisprojekt}

\subsection{Aufgabe: Multi-Agent Customer Support System}

Baue ein Multi-Agent-System für den Kundensupport.

\textbf{Anforderungen:}
\begin{itemize}[leftmargin=*]
    \item Implementiere drei spezialisierte Agenten: \texttt{OrderAgent} (Bestellstatus, Retoure), \texttt{KnowledgeAgent} (FAQ, Produktinfos), \texttt{EscalationAgent} (menschlicher Support bei komplexen Fällen)
    \item Ein \texttt{SupervisorAgent} analysiert die Anfrage und delegiert an den richtigen Worker
    \item Implementiere einen \texttt{SafetyGuard} mit max\_steps, Token-Budget und Tool-Whitelist
    \item Füge Human-in-the-Loop für Retoure-Freigaben hinzu
    \item Logge jeden Agent-Schritt (Thought, Action, Observation, Result)
    \item Implementiere ein Dashboard, das die Agent-Performance zeigt (Schritte pro Task, Erfolgsrate, Kosten)
\end{itemize}

\textbf{Testdaten:} 10 Support-Tickets (4 Bestellstatus, 3 FAQ, 3 komplexe Fälle).

\textbf{Erfolgskriterium:} 80\,\% der Tickets werden automatisiert gelöst. Alle kritischen Aktionen (Retoure > €100) landen im HITL. Kein Agent überschreitet max\_steps oder Budget.

% ===========================================================%
% INTERVIEWFRAGEN
% ===========================================================%
\section{Interviewfragen}

\begin{enumerate}[leftmargin=*]
    \item \textbf{Was ist der Unterschied zwischen einem RAG-System und einem AI Agenten?}

    \textbf{Antwort:} RAG reichert einen einzelnen Prompt mit Kontext an und generiert eine Antwort. Ein Agent führt einen Multi-Step-Loop aus: plant, ruft Tools auf, bewertet Ergebnisse und wiederholt. RAG ist passiv, Agent ist aktiv.

    \item \textbf{Erkläre den ReAct-Loop.}

    \textbf{Antwort:} Thought (LLM denkt nach) $\rightarrow$ Action (Tool-Aufruf) $\rightarrow$ Observation (Ergebnis) $\rightarrow$ Thought (nächster Schritt). Wiederholen, bis das Ziel erreicht ist. Dann: Answer.

    \item \textbf{Wann verwendest du Single-Agent, wann Multi-Agent?}

    \textbf{Antwort:} Single-Agent für klar definierte Aufgaben mit wenigen Tools (3--5). Multi-Agent bei komplexen Aufgaben, die verschiedene Expertise erfordern (Recherche + Analyse + Bericht). Multi-Agent ist teurer und langsamer, aber robuster.

    \item \textbf{Wie verhinderst du, dass ein Agent in einer Endlosschleife steckt?}

    \textbf{Antwort:} max\_steps (10--15), Timeout (30--60s), Token-Budget pro Run, Loop-Detection (wiederholtes Aufrufen des gleichen Tools mit gleichen Parametern). Bei Detektion: abbrechen und Fallback-Strategie ausführen.

    \item \textbf{Was ist Parallel Tool Calling und warum ist es wichtig?}

    \textbf{Antwort:} Das LLM kann mehrere Tools in einer Response aufrufen, statt nacheinander. Reduziert die Anzahl der LLM-Calls um 30--50\,\% und senkt die Latenz drastisch. OpenAI unterstützt es nativ, Anthropic nicht.

    \item \textbf{Wie schützt du einen Agenten vor schädlichen Tool-Ergebnissen?}

    \textbf{Antwort:} Safety-Guard prüft Tool-Output auf Injection-Patterns. Tool-Whitelist begrenzt, welche Tools aufgerufen werden dürfen. Output-Filter blockiert schädliche Anweisungen. HITL für kritische Tools.

    \item \textbf{Was ist der Unterschied zwischen ReAct und Plan-and-Solve?}

    \textbf{Antwort:} ReAct plant von Schritt zu Schritt (reaktiv, flexibel). Plan-and-Solve erstellt zuerst einen vollständigen Plan und führt ihn dann aus (strukturiert, vorhersagbar). Plan-and-Solve ist besser für komplexe Aufgaben mit vielen Abhängigkeiten.

    \item \textbf{Wie implementierst du Human-in-the-Loop?}

    \textbf{Antwort:} Approval-System mit pending\_approvals-Queue. Bei kritischen Aktionen (Zahlung, Löschung) pausiert der Agent und sendet eine Benachrichtigung (Slack, E-Mail). Der Mensch approved oder rejected. Der Agent wartet auf die Entscheidung (mit Timeout).

    \item \textbf{Wie evaluierst du die Qualität eines Agenten?}

    \textbf{Antwort:} Task-Success-Rate (wurde das Ziel erreicht?), Steps-per-Task (Effizienz), Cost-per-Task (Kosten), Tool-Selection-Accuracy (wurde das richtige Tool gewählt?), Error-Rate (wie oft ist ein Schritt fehlgeschlagen?).

    \item \textbf{Dein Agent ruft immer wieder das falsche Tool auf. Was machst du?}

    \textbf{Antwort:} (1) Tool-Descriptions prüfen und verbessern. (2) Tool-Namen präziser machen. (3) Beispiel-Interaktionen im System-Prompt einbauen. (4) Tool-Selection-Logging aktivieren, um zu sehen, warum das Modell das falsche Tool wählt. (5) Anzahl der Tools reduzieren.
\end{enumerate}

% ===========================================================%
% WEITERFÜHRENDE RESSOURCEN
% ===========================================================%
\section{Weiterführende Ressourcen}

\subsection{Dokumentationen}

\begin{itemize}[leftmargin=*]
    \item LangGraph Docs: \href{https://langchain-ai.github.io/langgraph/}{langchain-ai.github.io/langgraph/}
    \item OpenAI Function Calling: \href{https://platform.openai.com/docs/guides/function-calling}{platform.openai.com/docs/guides/function-calling}
    \item Anthropic Tool Use: \href{https://docs.anthropic.com/en/docs/build-with-claude/tool-use}{docs.anthropic.com/en/docs/build-with-claude/tool-use}
    \item CrewAI Docs: \href{https://docs.crewai.com}{docs.crewai.com}
    \item AutoGen Docs: \href{https://microsoft.github.io/autogen/}{microsoft.github.io/autogen/}
\end{itemize}

\subsection{Tools}

\begin{itemize}[leftmargin=*]
    \item \textbf{LangGraph}: Graph-basierte Agent-Orchestrierung
    \item \textbf{CrewAI}: Rollenbasierte Multi-Agent-Systeme
    \item \textbf{AutoGen (Microsoft)}: Multi-Talk-Messaging zwischen Agenten
    \item \textbf{OpenAI Assistants API}: Managed Agent Platform (serverless)
    \item \textbf{Braintrust}: \href{https://www.braintrust.dev}{braintrust.dev} -- Agent-Evaluierung und Observability
\end{itemize}

\subsection{Literatur}

\begin{itemize}[leftmargin=*]
    \item Yao et al. (2023): "ReAct: Synergizing Reasoning and Acting in Language Models"
    \item Wang et al. (2023): "Plan-and-Solve Prompting: Improving Zero-Shot Chain-of-Thought Reasoning"
    \item Xi et al. (2023): "The Rise and Potential of Large Language Model Based Agents: A Survey"
    \item Weng (2024): "LLM Agents" -- \href{https://lilianweng.github.io/posts/2023-06-23-agent/}{lilianweng.github.io}
\end{itemize}

\textbf{Nächster Schritt:} Mit dem Agent-Wissen geht es zur Evaluation -- wie du die Qualität deiner LLM-Applikationen misst und verbesserst.
